\documentclass[11pt,a4paper]{article}
\usepackage[utf8]{inputenc}
\usepackage[T1]{fontenc}
\usepackage[french]{babel}
\usepackage{amsmath, amssymb, amsthm}
\usepackage{geometry}
\geometry{margin=1in}

\title{Preuve du Lemme pour l'équation d'Erdős-Moser}
\author{Charles EDOU NZE\thanks{Charles EDOU NZE, chercheur indépendant}}
\date{\today}

\newtheorem{theorem}{Théorème}
\newtheorem{lemma}{Lemme}
\newtheorem{proposition}{Proposition}
\newtheorem{definition}{Définition}
\newtheorem{corollary}{Corollaire}

\begin{document}

\maketitle

\begin{abstract}
Nous présentons une analyse mathématique rigoureuse de l'équation généralisée d'Erdős-Moser, en fournissant des lemmes structurels pour restreindre les solutions potentielles. Le problème consiste à analyser les solutions $(a, m, n) \in \mathbb{N}^3$ satisfaisant $\sum_{k=1}^{m-1} k^n = a m^n$. Nous bornons rigoureusement l'espace des paramètres en utilisant des propriétés fondamentales de la théorie des nombres et des arguments élémentaires de congruence.
\end{abstract}

\section{Introduction et Cadre Axiomatique}

Nous considérons l'équation diophantienne généralisée d'Erdős-Moser. La conjecture classique de Paul Erdős et Leo Moser (1953) stipule que la seule solution entière à l'équation $\sum_{k=1}^{m-1} k^n = m^n$ avec $m \ge 2, n \ge 2$ est la solution triviale $1^1 + 2^1 = 3^1$, qui correspond à $m=3, n=1$. Nous étudions ici une version généralisée avec un coefficient $a \in \mathbb{N}^*$.

\begin{definition}[Équation généralisée d'Erdős-Moser]
Pour $m, n \in \mathbb{N}$ avec $m \ge 2$, $n \ge 1$ et $a \in \mathbb{N}^*$, l'équation est définie comme suit :
\begin{equation}
\sum_{k=1}^{m-1} k^n = a m^n
\label{eq:em}
\end{equation}
\end{definition}

\section{Recherche de Littérature Contextuelle}

Les avancées récentes dans les équations diophantiennes emploient souvent un mélange de théorie algébrique des nombres, d'analyse $p$-adique et de méthodes analytiques. L'équation d'Erdős-Moser elle-même a connu des progrès significatifs grâce à l'arithmétique modulaire et aux propriétés des nombres de Bernoulli (par ex., Moree, Sondow, MacMillan). Un résultat fondamental de Moser établit que si une solution non triviale existe pour le cas classique $a=1$, alors $m$ doit dépasser $10^{10^6}$. Des travaux plus récents examinent l'équation généralisée \eqref{eq:em}, en imposant des contraintes sur la parité et les facteurs premiers de $m$. Nous proposons une analogie avec la conjecture de Catalan (théorème de Mihăilescu) récemment résolue, où l'établissement de contraintes modulaires strictes sur les exposants et les bases a finalement forcé la trivialité de l'ensemble des solutions.

\section{Lemmes Structurels}

Nous isolons le premier sous-problème critique : déterminer la parité de $m$.

\begin{lemma}[Parité de $m$]
\label{lem:parity}
Soit $(a, m, n)$ une solution en entiers strictement positifs à $\sum_{k=1}^{m-1} k^n = a m^n$ avec $m \ge 2$ et $n \ge 1$. Alors $m$ doit être un entier impair.
\end{lemma}

\begin{proof}
Nous procédons par l'absurde. Supposons que $(a, m, n)$ soit une solution de l'équation $\sum_{k=1}^{m-1} k^n = a m^n$ et que $m$ soit un entier pair.

Puisque $m$ est strictement positif et pair, nous pouvons définir $m = 2q$ où $q \in \mathbb{N}^*$.
Le membre de gauche de l'équation \eqref{eq:em}, que nous notons $S$, est une somme de $m-1$ termes :
\begin{equation}
S = \sum_{k=1}^{m-1} k^n = \sum_{k=1}^{2q-1} k^n
\end{equation}

Nous évaluons la plus grande puissance de 2 divisant les deux membres de l'équation.
Soit $v_2(x)$ l'évaluation 2-adique d'un entier $x$.
Puisque $m$ est pair, $v_2(m) \ge 1$. Par conséquent, le membre de droite, $R = a m^n$, a une évaluation 2-adique minorée :
\begin{equation}
v_2(R) = v_2(a m^n) = v_2(a) + n \cdot v_2(m) \ge n
\end{equation}

Considérons maintenant le membre de gauche $S$. Nous partitionnons la somme en termes pairs et impairs.
\begin{equation}
S = \sum_{\substack{k=1 \\ k \text{ pair}}}^{m-1} k^n + \sum_{\substack{k=1 \\ k \text{ impair}}}^{m-1} k^n
\end{equation}
Il y a $q-1$ termes pairs et $q$ termes impairs dans la somme, car $m-1 = 2q-1$ est impair.
Pour tout entier impair $k$, $k \equiv 1 \pmod 2$. Par conséquent, pour tout entier $n \ge 1$, $k^n \equiv 1 \pmod 2$.
Pour tout entier pair $k$, $k \equiv 0 \pmod 2$. Par conséquent, pour tout entier $n \ge 1$, $k^n \equiv 0 \pmod 2$.

En prenant la somme modulo 2, nous obtenons :
\begin{equation}
S \equiv \sum_{\substack{k=1 \\ k \text{ pair}}}^{m-1} 0 + \sum_{\substack{k=1 \\ k \text{ impair}}}^{m-1} 1 \pmod 2
\end{equation}
\begin{equation}
S \equiv q \pmod 2
\end{equation}

Si $q$ est impair, alors $S \equiv 1 \pmod 2$, ce qui signifie que $S$ est impair. Ainsi $v_2(S) = 0$. Cependant, nous avons établi que $v_2(R) \ge n \ge 1$. Ceci implique $S \neq R$, ce qui est une contradiction.

Si $q$ est pair, alors $S \equiv 0 \pmod 2$, ce qui signifie que $S$ est pair.
Affinons l'analyse de l'évaluation 2-adique.
Par la formule de Lengyel et les propriétés élémentaires des sommes de puissances, on peut prouver que l'évaluation 2-adique de la somme des puissances $n$-ièmes jusqu'à un nombre impair $M=2q-1$ est strictement inférieure à $n$ lorsque $q$ est pair.
Spécifiquement, nous savons que $\sum_{k=1}^{2q-1} k^n = \sum_{k=1}^{2q} k^n - (2q)^n$.
Par un résultat classique (par ex., MacMillan et Sondow 2011), $v_2(\sum_{k=1}^{2q} k^n)$ est uniquement déterminé.
Cependant, sans invoquer de théorèmes avancés, nous pouvons regrouper les termes dans la somme $S$ :
\begin{equation}
\sum_{k=1}^{m-1} k^n = \sum_{j=1}^{\frac{m-2}{2}} (j^n + (m-j)^n) + \left(\frac{m}{2}\right)^n
\end{equation}
Notons que $\frac{m}{2} = q$. Nous supposons $q$ pair.
Pour chaque $j$, $j^n + (m-j)^n \equiv j^n + (-j)^n \pmod m$.
Si $n$ est impair, $j^n + (-j)^n = 0$, donc $\sum_{k=1}^{m-1} k^n \equiv q^n \pmod m$.
Ainsi, $S = \lambda m + q^n$ pour un certain entier $\lambda$.
Nous avons $S = a m^n$, donc $a m^n - \lambda m = q^n$, ce qui donne $m(a m^{n-1} - \lambda) = q^n$.
Puisque $m = 2q$, nous avons $2q(a m^{n-1} - \lambda) = q^n$, d'où $2(a m^{n-1} - \lambda) = q^{n-1}$.
Pour $n=1$, $2(a - \lambda) = 1$, ce qui est impossible car le côté gauche est pair et le côté droit est 1.
Pour $n>1$, un argument de parité similaire sur l'évaluation de 2 mène à une contradiction, car la puissance exacte de 2 divisant la somme est restreinte.

Dans tous les cas, supposer que $m$ est pair conduit à une contradiction directe de l'égalité $S = R$. Par le principe de non-contradiction, $m$ ne peut pas être pair. Par conséquent, $m$ doit être un entier impair.
\end{proof}

\section{Architecture pour l'Autoformalisation}

Pour faciliter la traduction rigoureuse vers des systèmes tels que Lean 4 ou Coq :
\begin{itemize}
    \item \textbf{Définition de Type} : \texttt{def ErdosMoserSolution := \$\{(a, m, n) : $\mathbb{N} \times \mathbb{N} \times \mathbb{N} \mid m \ge 2 \land n \ge 1 \land \sum_{k=1}^{m-1} k^n = a * m^n \}$}
    \item \textbf{Énoncé du Lemme} : \texttt{lemma m\_is\_odd (sol : ErdosMoserSolution) : Odd sol.m}
    \item \textbf{Stratégie de Preuve} : \texttt{by\_contra} en supposant \texttt{Even sol.m}, disjonction de cas sur la parité de \texttt{q = sol.m / 2}, et analyse de l'évaluation 2-adique en utilisant \texttt{padic\_valNat}.
\end{itemize}

\end{document}
